%%%%%%%%%%%%%%%%%%%%%%%%%%%%%%%%%%%%%%%
% Cover Letter - AeroVironment
% Software Engineering Intern (Embedded Focus), Simi Valley CA
% Plain XeLaTeX — matches simple resume style (Spectral font)
%%%%%%%%%%%%%%%%%%%%%%%%%%%%%%%%%%%%%%%

\documentclass[11pt]{article}

\usepackage[margin=0.75in, top=0.7in, bottom=0.7in]{geometry}
\usepackage[hidelinks]{hyperref}
\usepackage{parskip}
\usepackage{fontspec}

\setmainfont{Spectral}[
  Path = ../../cv/fonts/,
  UprightFont    = Spectral-Regular.ttf,
  BoldFont       = Spectral-Bold.ttf,
  ItalicFont     = Spectral-Italic.ttf,
  BoldItalicFont = Spectral-BoldItalic.ttf
]

\setlength{\parindent}{0pt}
\setlength{\parskip}{14pt}
\linespread{1.15}
\raggedright

% No hyphenation
\hyphenpenalty=10000
\exhyphenpenalty=10000

\pagestyle{empty}
\urlstyle{same}

\begin{document}

% ── DATE ─────────────────────────────────────────────────────────────────────
\begin{flushright}
  \today
\end{flushright}

% ── BODY ─────────────────────────────────────────────────────────────────────
Dear AV Hiring Team,

I'm applying for the Software Engineering Intern (Embedded Focus) role in
Simi Valley for Summer 2027. I'm a Computer Science student at RIT in the
combined bachelor's and master's program, with an expected graduation date
of December 2027, and have a strong interest in firmware and embedded
software. At GE Vernova, I built a full I/O subsystem from scratch in C on
an STM32F7 for industrial radios. On the RIT Electric Vehicle Team, I'm the
vehicle control unit (VCU) Software Lead, building real-time ThreadX
firmware for our track motorcycle. In my free time I'm building an
autonomous fixed-wing flight controller on an ESP32S3, with PID
stabilization, sensor fusion, and GPS waypoint following validated across
an 80-test SITL regression suite; the full airframe recently completed its
first flight in manual mode.

What got me interested in AeroVironment specifically is how tightly a
platform like Switchblade's size and power budget are tuned. The 300
variant's full launch-ready weight is under eight pounds, and flight
endurance changes directly depending on whether it's carrying a warhead or
just sensors. That's the same kind of power, weight, and compute tradeoff I
deal with on my own flight controller, where every component on the board
competes with flight time. I want to help build the embedded software that
works inside those constraints.

Sensor and real-time control firmware is where I spend most of my time.
I've written interrupt-driven I/O and real-time control loops for both
bare-metal and RTOS targets, in ThreadX and FreeRTOS, and I debug at the
hardware level regularly with a Saleae logic analyzer, an oscilloscope, and
a J-Link Segger. My flight controller runs PID stabilization and GPS
waypoint navigation on an ESP32S3 over SPI, I2C, and UART sensor buses,
validated in simulation before any of it flew. I've also worked across SPI,
I2C, UART, ADC, and PWM building both the GE Vernova I/O subsystem and
EVT's hardware-agnostic peripheral library, writing drivers directly from
microcontroller datasheets and reference manuals. At Lockheed Martin, I
built a code-review automation tool that became a mandatory part of the
team's merge request process, so I'm just as comfortable with the
peer-review and configuration-management side of embedded development.

I would welcome the opportunity to talk more about where I could
contribute. Thank you for your consideration.

Thanks,

Travis Brown

\end{document}
